\documentclass[a4paper,12pt]{article}

\usepackage[a4paper,margin=16mm]{geometry}
\usepackage[german,main=french]{babel}
\usepackage{fontspec}
\usepackage{microtype}

\setmainfont{EBGaramond}
[
  Extension      = .otf,
  UprightFont    = *-Regular,
  ItalicFont     = *-Italic,
  BoldFont       = *-Bold,
  BoldItalicFont = *-BoldItalic,
  Numbers        = Lowercase
]

\usepackage{paracol}
\newcommand\chunks[2]{%
\begin{leftcolumn*}\begin{otherlanguage}{german}%
{#1}%
\end{otherlanguage}\end{leftcolumn*}%
\begin{rightcolumn}\begin{otherlanguage}{french}%
{#2}%
\end{otherlanguage}\end{rightcolumn}%
}

\newcommand\strophe[1]{\subsection*{\normalsize\normalfont\itshape #1}}

\begin{document}\thispagestyle{empty}
\setlength{\columnsep}{8mm}


\begin{paracol}{2}

  \chunks{\section*{Der Vogelfänger bin ich ja}}
         {\section*{L’attrapeur d’oiseaux, oui c’est moi}}

  \chunks{\noindent\raggedright\small Wolfgang Amadeus Mozart, \textit{Die Zauberflöte}}
         {\noindent\raggedright\small Wolfgang Amadeus Mozart, \textit{La Flûte enchantée}}

  \chunks{\noindent\raggedright\small Libretto von Emanuel Schikaneder (1791)}
         {\noindent\raggedright\small Livret d’Emanuel Schikaneder (1791)}

  \chunks{}
         {\noindent\raggedright\small Texte français : Jean Spenlehauer}

  \chunks{\strophe{Erster Aufzug, zweiter Auftritt, Nr.~2, Arie}}
         {\strophe{Premier acte, deuxième scène, n\textsuperscript{o}~2, air}}

  \chunks{\noindent\small(Papageno kommt den Fußsteig herunter, hat auf dem Rücken eine große Vogelsteige, die hoch über den Kopf geht, worin verschiedene Vögel sind; auch hält er mit beiden Händen ein Faunen-Flötchen, pfeift und singt.)}
         {\noindent\small(Papageno arrive, descendant le sentier, ayant sur le dos une grande cage dépassant largement sa tête et où se trouvent différents oiseaux ; il a aussi une flûte de Pan, il siffle et chante.)}

  \chunks{\begin{verse}
          \textsc{Papageno}\\[\medskipamount]
          Der Vogelfänger bin ich ja,\\
          Stets lustig, heißa! hopsasa!\\
          Ich Vogelfänger bin bekannt\\
          Bei Alt und Jung im ganzen Land.\\
          Weiß mit dem Locken umzugehn,\\
          Und mich aufs Pfeifen zu verstehn.\\
          Drum kann ich froh und lustig sein,\\
          Denn alle Vögel sind ja mein.
          \end{verse}}
         {\begin{verse}
          \textsc{Papageno}\\[\medskipamount]
          L’attrapeur d’oiseaux, oui c’est moi\\
          Toujours gai, heissa, hopsasa !\\
          Moi, attrapeur d’oiseau, je suis connu\\
          Par grands et petits dans tout le pays.\\
          Je m’y connais en appeaux\\
          Et en sifflets.\\
          Donc je peux être joyeux et gai\\
          Car tous les oiseaux sont à moi.
          \end{verse}}

  \chunks{\begin{verse}
          Der Vogelfänger bin ich ja,\\
          Stets lustig, heißa! hopsasa!\\
          Ich Vogelfänger bin bekannt,\\
          Bei Alt und Jung im ganzen Land.\\
          Ein Netz für Mädchen möchte ich,\\
          Ich fing’ sie dutzendweiß für mich.\\
          Dann sperrte ich sie bei mir ein,\\
          Und alle Mädchen wären mein.
          \end{verse}}
         {\begin{verse}
          L’attrapeur d’oiseaux, oui c’est moi\\
          Toujours gai, heissa, hopsasa !\\
          Moi, attrapeur d’oiseau, je suis connu\\
          Par grands et petits dans tout le pays.\\
          J’aimerais avoir un filet à filles\\
          Je m’en attraperais par douzaines\\
          Puis je les enfermerais chez moi\\
          Et toutes les filles seraient à moi.
          \end{verse}}

  \chunks{\begin{verse}
          Wenn alle Mädchen wären mein,\\
          So tauschte ich brav Zucker ein:\\
          Die, welche mir am liebsten wär’,\\
          Der gäb’ ich gleich den Zucker her.\\
          Und küßte sie mich zärtlich dann,\\
          Wär’ sie mein Weib und ich ihr Mann.\\
          Sie schlief’ an meiner Seite ein,\\
          Ich wiegte wie ein Kind sie ein.
          \end{verse}}
         {\begin{verse}
          Si toutes les filles étaient à moi,\\
          Je les échangerais contre du sucre :\\
          A celle que j’aimerais le plus,\\
          Je donnerais ce sucre.\\
          Et si elle m’embrassait tendrement,\\
          Elle serait ma femme et moi son mari.\\
          Elle s’endormirait à mes côtés,\\
          Je la bercerais comme un enfant.
          \end{verse}}
\end{paracol}
\end{document}
